% ============================================================================
% COURS 03 — LOIS ENTRÉE-SORTIE — PARTIE A
% Source : 3-LoiES20222.docx
% ============================================================================

\section{Lois entrée-sortie : objectifs et mises en situation}
\label{sec:les-objectifs}

\begin{TLThreeMethod}{Objectifs du cours}
Donner les outils permettant de déterminer la loi entrée-sortie d'un
mécanisme, en position puis en vitesse, à partir de son modèle
cinématique.
\end{TLThreeMethod}

\begin{center}
\small
\begin{tabularx}{\linewidth}{>{\bfseries}p{2cm}X c}
\toprule
Je connais & Connaissances attendues & \\
\midrule
& La démarche de détermination d'une loi entrée-sortie. & $\square$\\
& Les notions de modèle géométrique direct et inverse. & $\square$\\
& Les coordonnées articulaires et opérationnelles. & $\square$\\
& La relation de Chasles, les opérations vectorielles et le produit scalaire. & $\square$\\
\midrule
Je sais & Identifier les paramètres d'entrée, de sortie et les paramètres intermédiaires. & $\square$\\
& Écrire et projeter une fermeture géométrique. & $\square$\\
& Éliminer un angle ou une distance intermédiaire. & $\square$\\
& Déduire une loi en vitesse d'une loi en position. & $\square$\\
\bottomrule
\end{tabularx}
\end{center}

\subsection{Cas n°1 : dimensionnement d'une course de vérin}

Une E.P.A.S. est une échelle pivotante automatique à commande
séquentielle montée sur un camion de pompiers. Le vérin de dressage doit
faire passer l'échelle de la position basse, $\theta=10^\circ$, à la
position haute, $\theta=50^\circ$.

\begin{center}
\begin{minipage}[c]{.44\linewidth}
\TLThreeImage[\linewidth]{images/image5.png}
\end{minipage}\hfill
\begin{minipage}[c]{.52\linewidth}
\TLEPASGeometry
\end{minipage}
\end{center}

\begin{TLThreeApp}{Problème posé}
Déterminer la relation $r=f(\theta)$ entre la longueur du vérin et
l'angle de dressage, puis en déduire la course minimale du vérin entre
les deux positions extrêmes.
\end{TLThreeApp}

\subsection{Cas n°2 : espace de travail d'un robot}

Le robot porte-instrument aide le praticien à positionner une aiguille,
une sonde ou un endoscope. Pour garantir la précision de l'intervention,
il faut connaître son \textbf{espace de travail}, c'est-à-dire l'ensemble
des positions atteignables par l'effecteur.

\begin{center}
\TLThreeImage[.36\linewidth]{images/image8.png}\hfill
\TLThreeImage[.48\linewidth]{images/image9.png}
\end{center}

La question posée est différente de celle d'une chaîne fermée : les
paramètres articulaires étant connus, il faut calculer les coordonnées
de l'outil dans le repère du bâti.

\subsection{Cas n°3 : cylindrée d'un micromoteur}

Un micromoteur d'avion de modélisme transforme la rotation du vilebrequin
en translation alternative du piston par un mécanisme bielle-manivelle.

\begin{center}
\TLThreeImage[.48\linewidth]{images/image10.png}\hfill
\TLThreeImage[.44\linewidth]{images/image11.png}
\end{center}

\TLBielleManivelleGeometry

Le moteur ENIA--32K-C possède un piston de diamètre $24\ \mathrm{mm}$ et
une cylindrée annoncée de $10\ \mathrm{cm^3}$.

\begin{TLThreeApp}{Questions de conception}
\begin{enumerate}
  \item Identifier le mouvement d'entrée et le mouvement de sortie.
  \item Établir la loi reliant l'angle du vilebrequin à la position du piston.
  \item Déterminer la course puis vérifier la cylindrée annoncée.
\end{enumerate}
\end{TLThreeApp}

\section{Structures fermées et paramètres de mouvement}
\label{sec:les-fermees}

Un mécanisme peut transformer un mouvement et des efforts pour les
adapter au besoin. Lorsqu'il comporte une chaîne cinématique fermée, la
mise en mouvement d'une liaison entraîne tous les solides de la boucle.

\begin{TLThreeDef}{Loi entrée-sortie en position}
La loi entrée-sortie cinématique en position est la relation
mathématique entre le paramètre de position de l'entrée et celui de la
sortie :
\[
q_s=f(q_e).
\]
Elle ne doit plus contenir les paramètres associés aux liaisons
intermédiaires.
\end{TLThreeDef}

Pour un bielle-manivelle utilisé comme compresseur, l'entrée est souvent
l'angle $\alpha$ et la sortie la position $x$ du piston. Dans un moteur
thermique, le sens du transfert de puissance est inversé et $x$ peut
être considéré comme l'entrée.

\begin{TLThreeDef}{Course, cylindrée et homocinétisme}
\textbf{Course :} distance entre les deux positions extrêmes d'un
coulisseau ou d'un piston.

\textbf{Cylindrée :} volume balayé par le piston pendant sa course,
multiplié par le nombre de cylindres :
\[
V_c=n\,\frac{\pi D^2}{4}\,C.
\]

\textbf{Mécanisme homocinétique :} mécanisme dont le rapport des vitesses
est constant et égal à l'unité. Un cardan simple n'est pas
homocinétique ; deux cardans correctement montés peuvent l'être.
\end{TLThreeDef}

\section{Structures ouvertes : coordonnées articulaires et opérationnelles}
\label{sec:les-ouvertes}

Dans une chaîne ouverte de type robot série :
\begin{itemize}
  \item les \textbf{coordonnées articulaires} sont les paramètres de
  position des liaisons : angles de pivots et courses de glissières ;
  \item les \textbf{coordonnées opérationnelles} décrivent la position
  et éventuellement l'orientation de l'effecteur dans le repère de
  référence.
\end{itemize}

Pour un robot plan à quatre degrés de liberté, on peut écrire :
\[
\mathbf q=(\alpha,\beta,\gamma,\delta),
\qquad
\mathbf X=(x_P,y_P,z_P).
\]

\begin{TLThreeDef}{Modèles direct et inverse}
Le \textbf{modèle géométrique direct} calcule les coordonnées
opérationnelles à partir des coordonnées articulaires :
\[
\mathbf X=F(\mathbf q).
\]
Le \textbf{modèle géométrique inverse} cherche les coordonnées
articulaires permettant d'atteindre une position imposée :
\[
\mathbf q=F^{-1}(\mathbf X).
\]
Il peut posséder plusieurs solutions, une solution unique ou aucune
solution.
\end{TLThreeDef}

\section{Rappels sur les vecteurs}
\label{sec:les-vecteurs}

\subsection{Composantes et norme}

Dans une base orthonormée directe
$\mathcal B=(\vec i,\vec j,\vec k)$ :
\[
\vec V=X\vec i+Y\vec j+Z\vec k
\quad\Longleftrightarrow\quad
[\vec V]_{\mathcal B}=\begin{pmatrix}X\\Y\\Z\end{pmatrix}.
\]
Sa norme vaut :
\[
\boxed{\|\vec V\|=\sqrt{X^2+Y^2+Z^2}.}
\]

\subsection{Addition, multiplication et relation de Chasles}

\[
\vec U+\vec V=
(U_x+V_x)\vec i+(U_y+V_y)\vec j+(U_z+V_z)\vec k,
\]
\[
\lambda\vec U=(\lambda U_x)\vec i+(\lambda U_y)\vec j+(\lambda U_z)\vec k.
\]

Pour trois points $A$, $B$ et $C$ :
\[
\boxed{\overrightarrow{AB}+\overrightarrow{BC}=\overrightarrow{AC}.}
\]
Dans une chaîne fermée $A_1A_2\ldots A_nA_1$ :
\[
\boxed{\overrightarrow{A_1A_2}+\overrightarrow{A_2A_3}+
\cdots+\overrightarrow{A_nA_1}=\vec0.}
\]

\begin{TLThreeApp}{Vitesse d'un bateau}
Une embarcation se déplace à $40\ \mathrm{km\,h^{-1}}$ par rapport à
l'eau, tandis que le courant vaut $15\ \mathrm{km\,h^{-1}}$. La vitesse
par rapport à la rive est la somme vectorielle des deux vitesses.
Construire la somme et déterminer sa norme pour une direction du bateau
faisant $50^\circ$ avec la rive.
\end{TLThreeApp}

\subsection{Produit scalaire}

\begin{TLThreeDef}{Définition et expression analytique}
\[
\vec U\cdot\vec V=\|\vec U\|\,\|\vec V\|\cos(\vec U,\vec V).
\]
Dans une base orthonormée :
\[
\boxed{\vec U\cdot\vec V=U_xV_x+U_yV_y+U_zV_z.}
\]
Le produit scalaire est nul si l'un des vecteurs est nul ou si les deux
vecteurs sont orthogonaux.
\end{TLThreeDef}

Les propriétés utiles sont la symétrie, la distributivité et la
compatibilité avec la multiplication par un réel :
\[
\vec U\cdot\vec V=\vec V\cdot\vec U,
\qquad
\vec U\cdot(\vec V+\vec W)=\vec U\cdot\vec V+\vec U\cdot\vec W.
\]

\begin{TLThreeApp}{Calcul direct}
Avec
$\vec V_1=2\vec x+2\vec y-3\vec z$ et
$\vec V_2=-\vec x+6\vec z$, calculer
$\vec V_1\cdot\vec V_2$ puis l'angle entre les deux vecteurs.
\end{TLThreeApp}
